\documentclass[12pt,a4paper]{article}
\usepackage{amsmath}
\usepackage{amssymb}
\usepackage{graphicx}
\usepackage{subcaption}
\usepackage{enumitem}
\usepackage{listings}
\usepackage{hyperref}
\usepackage{ulem}
\usepackage[table]{xcolor}
\usepackage{amsthm}
\usepackage{tcolorbox}
\usepackage{array}
\usepackage{mdframed}
\usepackage{cancel}
\usepackage[margin=2.5cm]{geometry}

\setlength{\parindent}{0pt}
\setlength{\parskip}{0pt}
\setcounter{MaxMatrixCols}{11}

\begin{document}
\begin{titlepage}
\centering
\vspace*{\fill}
{\Huge Algebra}\\[0.5cm]
{\Huge Home Work 9}\\[1.5cm]
{\Large Student: 26-712-224}\\[0.5cm]
{\Large \today}\par
\vspace{2cm}
{\normalsize ``I wrote these solutions fully myself und would like to receive feedback on these''\par}
\vspace*{\fill}
\end{titlepage}
\noindent\textbf{A)} For each of the following permutations find the decomposition into cycles with disjoint supports, then determine the order and the signature.
\begin{enumerate}
\renewcommand{\labelenumi}{(\arabic{enumi})}
\item $\sigma=\begin{pmatrix}
1&2&3&4&5&6&7&8&9\\
3&4&9&7&6&5&1&2&8
\end{pmatrix}$.
\item $\tau=\begin{pmatrix}
1&2&3&4&5&6&7&8&9\\
2&4&6&8&1&3&5&7&9
\end{pmatrix}$.
\item $\sigma\circ\tau^{-1}$.
\item $\sigma^2\circ\tau^3$.
\item $\begin{pmatrix}
1&2&\cdots&n&n+1&n+2&\cdots&2n\\
n+1&n+2&\cdots&2n&1&2&\cdots&n
\end{pmatrix}$.
\end{enumerate}

\vspace{1cm}
\noindent
Answer:\\
\\
(1) $\sigma=(1 \; 3 \; 9 \; 8 \; 2 \; 4 \; 7)(5 \; 6)$.  From which:
\begin{align*}
  \text{sgn}(\sigma) &= (-1)^6(-1)^1=-1 \\
  \text{ord}(\sigma) &= \text{lcm}(7,2)=14
\end{align*}
(2) $\tau=(1 \; 2 \; 4 \; 8 \; 7 \; 5)(3 \; 6)$.  From which:
\begin{align*}
  \text{sgn}(\tau) &= (-1)^5(-1)^1=1 \\
  \text{ord}(\tau) &= \text{lcm}(6,2)=6
\end{align*}
(3) $\tau^{-1} = (5 \; 7 \; 8 \; 4 \; 2 \; 1)(3 \; 6)$ and hence:
\begin{align*}
  \sigma \circ \tau^{-1} &= (1 \; 3 \; 9 \; 8 \; 2 \; 4 \; 7)(5 \; 6)(5 \; 7 \; 8 \; 4 \; 2 \; 1)(3 \; 6) \\
   &= (1 \; 6 \; 9 \; 8 \; 7 \; 2 \; 3 \; 5) \\
   \\
  \implies \text{sgn}(\sigma \circ \tau^{-1}) &= \text{sgn}(\sigma)\text{sgn}(\tau^{-1}) = \text{sgn}(\sigma)\text{sgn}(\tau) = (-1)(1)=-1\\
  \implies \text{ord}(\sigma \circ \tau^{-1}) &= 8 \;\;\;\;\;\;\;\;\; \text{(i.e. the length of the cycle)}
\end{align*}
(4) 
\begin{align*}
  \sigma^2 &= (1 \; 3 \; 9 \; 8 \; 2 \; 4 \; 7)^2 \\
   &= (1 \; 9 \; 2 \; 7 \; 3 \; 8 \; 4) \\
  \tau^3 &= (1 \; 2 \; 4 \; 8 \; 7 \; 5)^3(3 \; 6) \\
  &= (1 \; 8)(2 \; 7)(4 \; 5)(3 \; 6)
\end{align*}
From which:
\begin{align*}
  \sigma^2 \circ \tau^3 &= (1 \; 9 \; 2 \; 7 \; 3 \; 8 \; 4)(1 \; 8)(2 \; 7)(4 \; 5)(3 \; 6) \\
  &= (1 \; 4 \; 5)(2 \; 3 \; 6 \; 8 \; 9) \\
  \\
  \implies \text{sgn}(\sigma^2 \circ \tau^3) &= \text{sgn}(\sigma)^2 \text{sgn}(\tau)^3 = (-1)^2(1)^3=1\\
  \implies \text{ord}(\sigma^2 \circ \tau^3) &= \text{lcm}(3, 5)=15
\end{align*}

\newpage
\noindent
(5) Let:
\begin{align*}
  \sigma = \begin{pmatrix}
1&2&\cdots&n&n+1&n+2&\cdots&2n\\
n+1&n+2&\cdots&2n&1&2&\cdots&n
\end{pmatrix}
\end{align*}
Then we can also represent $\sigma$ as disjoint cycles:
\begin{align*}
  \sigma = \prod_{i=1}^n (i \; (n+i))
\end{align*}
So, it is clear that:
\begin{align*}
  \text{sgn}(\sigma) &= (-1)^n \\
  \text{ord}(\sigma) &= 2 \;\;\;\;\;\;\;\;\;\;\;\;\;\;\;\;\;\; \text{(as 2 is order for all the transpositions)}
\end{align*}

\newpage
\noindent\textbf{B)} Let $n\geq 3$ be an integer. The goal of this exercise is to show that the alternating group $A_n$ is generated by cycles of length $3$. To do so, we let $H\subset A_n$ be the subgroup generated by all the cycles $(uvw)$ where $u$, $v$ and $w$ are pairwise distinct integers in $\{1,\ldots,n\}$. The goal is to show that $H=A_n$.
\begin{enumerate}
\renewcommand{\labelenumi}{(\arabic{enumi})}
\item Recall that every even permutation $\sigma\in A_n$ is a product of an even number of transpositions. Deduce that it is enough to show that $(ab)(cd)\in H$ for $a,b,c,d\in\{1,\ldots,n\}$ with $a\neq b$ and $c\neq d$.
\item What happens if $\{a,b\}=\{c,d\}$?
\item Assume that $\operatorname{card}(\{a,b\}\cap\{c,d\})=1$. Show that $(ab)(cd)$ is a cycle of length $3$.
\item Assume that $\{a,b\}\cap\{c,d\}=\emptyset$. Show that $(ab)(cd)=(abc)(bcd)$.
\item Use the previous questions to conclude that $H=A_n$.
\end{enumerate}

\vspace{1cm}
\noindent
Answer:\\
\\
We should note the two following identities:
\begin{align*}
  (a \; b \; c) &= (a \; b)(b \; c) \;\;\;\;\;\;\;\;\;\;\;\;\;\;\;\;\;\;\;\;\;\;\;\;\;\;\;\;\;\;\; a \ne b; b \ne c; a \ne c; \;\; a,b,c \in \{1,2,...,n\} \tag{B.1}  \\
  \\
  (a \; b)(c \; d) &= [(a \; b)(b \; c)]\;[(b \; c)(c \; d)] \;\;\;\;\;\;\;\;\;\;\;\; a \ne b; c \ne d; \;\; a,b,c,d \in \{1,2,...,n\}  \\
  &= (a \; b \; c)(b \; c \; d) \;\;\;\;\;\;\;\;\;\;\;\;\;\;\;\;\;\;\;\;\;\;\;\;\; \text{(using (B.1))} \tag{B.2}
\end{align*}
(1) We define: $H = \{(u \; v \; w) \mid u,v,w \in \{1,2,3,..,9\} \;\; \text{with $u,v,w$ pairwise distinct} \}$.  If we assume that $(a \; b)(c \; d) \in H$ for $a,b,c,d \in \{1,2,...,n\}$ and
$a \ne b; c \ne d$.  Then it follows that $(a \; b)(c \; d)$ is the product of 3-cycles then every even sequence of transpositions can be broken-up into groups of two transpositions and will also be in $H$.\\
In none of the sequence of transpositions will we find that $a=b$ or $c=d$.\\
\\
(2) If $\{a,b\} = \{c,d\}$, then it must be the case that $a=c$ and $d=b$ (swapping to $a=d$ and $b=c$ makes no difference) and we find that we
have two identical transpositions - which cancel to give the identity element.\\
\\
(3) In the case that $\operatorname{card}(\{a,b\}\cap\{c,d\})=1$ we have, WLOG, that $b=c$ and we have from (B.1) that: $(a \; b)(b \; c) = (a \; b \; c)$ - as required.\\
\\
(4) In the case $\{a,b\}\cap\{c,d\}=\emptyset$ we must have two disjoint transpositions and then (B.2) gives us that:
$(a \; b)(c \; d) = (a \; b \; c)(b \; c \; d)$ - as required.\\
\\
(5) All possible combinations of an even number of transpositions have now been accounted for.\\
We have (i) identical transpositions - which cancel; (ii) transpositions with only element in common and (iii) transpoistions which are disjoint.  In all cases they can be written as 
a product of 3-cycles.  This, with (1), is enough to show that: $A_n = H$.

\newpage
\noindent\textbf{C)} Consider in the following permutation $\sigma$ in the symmetric group $S_9$:
\begin{align*}
\sigma=\begin{pmatrix}
1&2&3&4&5&6&7&8&9\\
4&5&6&7&8&9&1&2&3
\end{pmatrix}.
\end{align*}
\begin{enumerate}
\renewcommand{\labelenumi}{(\arabic{enumi})}
\item Decompose $\sigma$ into cycles and compute its order and signature.
\item Deduce from (1) that there exists a partition $\{1,\ldots,9\}=I_1\cup I_2\cup I_3$ such that $\operatorname{card}(I_i)=3$ and $\sigma(I_i)=I_i$ for $1\leq i\leq 3$.
\end{enumerate}
\noindent Let $C=C_{S_9}(\sigma)$ be the centralizer of $\sigma$ (i.e., $C=\{\tau\in S_9\mid\tau\sigma=\sigma\tau\}$).
\begin{enumerate}
\renewcommand{\labelenumi}{(\arabic{enumi})}
\setcounter{enumi}{2}
\item For $\tau\in C$, show that $\tau(I_i)\in\{I_1,I_2,I_3\}$ for all $1\leq i\leq 3$.
\item Deduce from the previous question that there exists a morphism of groups $\rho\colon C\to S_3$ such that $\tau(I_i)=I_{\rho(\tau)(i)}$ for all $1\leq i\leq 3$.
\item Show that there exists a unique morphism of groups
\begin{align*}
\epsilon\colon\mathbb{Z}/3\mathbb{Z}\times\mathbb{Z}/3\mathbb{Z}\times\mathbb{Z}/3\mathbb{Z}\to C
\end{align*}
given by
\begin{align*}
(e_1+3\mathbb{Z},e_2+3\mathbb{Z},e_3+3\mathbb{Z})\mapsto\alpha_1^{e_1}\alpha_2^{e_2}\alpha_3^{e_3}
\end{align*}
where $\alpha_1=(1\,4\,7)$, $\alpha_2=(2\,5\,8)$ and $\alpha_3=(3\,6\,9)$ are cycles in $S_9$.
\item Show that $\epsilon$ induces an isomorphism between $\mathbb{Z}/3\mathbb{Z}\times\mathbb{Z}/3\mathbb{Z}\times\mathbb{Z}/3\mathbb{Z}$ and $\operatorname{Ker}(\rho)$.
\item Show that $\rho$ is surjective. What is the cardinality of $C$? What is the cardinality of the conjugacy class of $\sigma$?
\end{enumerate}

\vspace{1cm}
\noindent
Answer:\\
\\
(1) Spliting $\sigma$ into distinct cycles gives: $\sigma=(1 \; 4 \; 7)(2 \; 5 \; 8)(3 \; 6 \; 9)$.\\
From which: $\text{sgn}(\sigma)=(-1)^2(-1)^2(-1)^2=1$ and $\text{ord}(\sigma)=\text{lcm}(3,3,3)=3$.\\
\\
(2) Given that $\sigma=(1 \; 4 \; 7)(2 \; 5 \; 8)(3 \; 6 \; 9)$ we see that we have 3 sets of numbers:
\begin{align*}
  I_1 &= \{1,4,7 \} \\
  I_2 &= \{2,5,8 \} \\
  I_3 & =\{3,6,9 \}
\end{align*}
which are affected by $\sigma$ and "rotates" the numbers each time the permutation is appled (e.g. for $I_1$, $\sigma(1)=4, \sigma^2(1)=7), \sigma^3(1)=1$).  This means the permutation of the elements in each set $I_i$ stays in that set.\\
Hence: 
\begin{align*}
  I_1 \cup I_2 \cup I_3 &= \{1,2,3,...,9\} \\
  \text{card}(I_i) &= 3 \;\;\;\;\;\;\;\;\;\;\;\;\;\;\;\;\;\; i \in \{1,2,3\}\\
  \sigma(I_i) &= I_i \;\;\;\;\;\;\;\;\;\;\;\;\;\;\;\;\;\; i \in \{1,2,3\}
\end{align*}
(3) Consider the group:
\begin{align*}
  G &= \{\sigma^0, \sigma^1, \sigma^2 \} \\
    &= \{I, (1 \; 4 \; 7)(2 \; 5 \; 8)(3 \; 6 \; 9), (1 \; 7 \; 4)(2 \; 8 \; 5)(3 \; 9 \; 6) \}
\end{align*}
I.e. $|G|=3$.  If we consider $X=\{1,2,..,9\}$ then X is a G-set.  Now $\sigma$ does not leave any element of 
X fixed, and hence: $\text{Stab}_G(s) = \{I\}; \forall s \in S_9$.  Applying the "Orbit-Stabilizer" theorem we conclude that each element of
X under $G$ will have an orbit of 3 elements.  Orbits are disjoint and partition X.  There

\end{document}